\scanpage{102}\mspage{93}
distinct de $P_a$ est $\mathrel{|\circ|}$ \unsure{à|:} $P_a$), d'où résulte aussitôt (à cause de l'antiréflexivité de $\mathrel{|\circ|}$) que $\mathrm{supp}^{\circ}(P_a)=\{P_a\}$. Si $a$ n'est \deleted{\unsure{p}}\unsure{ni|pas} plus petit ni plus grand élément, soient
\[
 \deleted{\unsure{x<a}}\qquad x<a<y\qquad .
\]
Alors $\deleted{\unsure{\exists|Z}}S_{x,a}\in\mathrm{cosupp}^{\circ}(P_a)$, donc $X\mathrel{|\circ|}S_{x,a}$, \deleted{donc} $\delta X\geq a$ ou $\delta X\leq x$. Mais on ne peut avoir $\delta X\leq x$, car on aurait $X\mathrel{|\circ|}a$ i.e. $X\in\mathrm{cosupp}^{\circ}(P_a)$, ce qui contredit $X\in\mathrm{supp}^{\circ}(P_a)$ (vu l'\unsure{antiréflexivité|antiréflexité}) :
\[
 \mathrm{cosupp}^{\circ}(\Phi)\cap\mathrm{supp}^{\circ}(\deleted{\unsure{\Psi}}\Phi)=\varnothing,\quad\text{pour tte partie}
\]
$\Phi$ de $\unsure{M|\mathcal M}$). Donc $\delta X\geq a$. L'\deleted{\unsure{m\={m}|=}}argument symétrique, utilisant $S_{a,y}$ \unsure{ald|.ld} $S_{xa}$, donnera $\delta X\leq a$, d'où la conclusion.

\uline{Cor. 2} Raisonnement analogue ; \deleted{\unsure{on}} \inserted{On} \deleted{Si $a$ est un \unsure{p}} prouve successivement
\[
 \delta X\geq a,\quad\text{puis }\delta X\leq b
\]
en distinguant, pour la première, le cas où $a$ est un plus petit élément (\unsure{trivial}) et le cas où il ne l'est pas, donc $\exists\,x$
\[
 x<a<b
\]
et on aura $S_{x,a}\in\mathrm{cosupp}^{\circ}(S_{ab})$ donc $X\mathrel{|\circ|}S_{x,a}$ donc $\delta X\leq x$ ou $\delta X\geq a$, et le premier cas est exclu comme devant \unsure{entrainer|entraîner} $X\mathrel{|\circ|}\unsure{\widetilde S_{ab}|S_{ab}}$, \unsure{exclu,|exclu.}
% archive sheet number: 102
